% page 379 of the facsimile (image p-378 of the scan)
% block 162.6 x 242.0 mm ; left margin 39.4 mm
\pgc{17.780mm}{0.000mm}{
\l{\cel{54}371}
\l{}
\l{\sou{mezotinto}.Graburo por qua onu preparas, unesme, la grano sur}
\l{\cel{2}la tota planko, til obtenar nigro perfekta, ed ube, pos de-}
\l{\cel{3}kalquir la desegnuro sur la kupro-plako, onu aplastas plu o}
\l{\cel{3}min multe la grano per skrapilo por la loki klara, la mi-}
\l{\cel{3}nuanci, la mi-ombri. - DEFI.}
\l{}
\l{\sou{mezurar}. (\sou{trans., per}). Evaluar (quanto)per lua raporto ye quan-}
\l{\cel{3}to determinita, sama-speca, selektita kom komparo-termino.-}
\l{\cel{3}EFI.}
\l{}
\l{\sou{\textquotedbl{}mi\textquotedbl{}.}\cel{2}(\sou{muziko}). Noto triesma di la gamo diatonika naturala,}
\l{\cel{3}t.e. di la \textquotedbl{}do\textquotedbl{}-gamo.}
\l{}
\l{\sou{mi-}\cel{3}Prefixo di qua la senco esas \textquotedbl{}duime\textquotedbl{}.}
\l{}
\l{\sou{miasmo}. Emanajo des-salubra,nociva, de materii putranta. - DEFIR}
\l{}
\l{\sou{miaular}. (\sou{netrans.}) (\sou{aludante kato}). Agar la krio qua esas propr}
\l{\cel{3}ye lua speco. - DEFIS.}
\l{}
\l{\sou{micelo}. (\sou{biol.}) La aglomerajo maxim mikra ek molekuli, qua havas}
\l{\cel{3}la propraji fizikala di korpo.}
\l{}
\l{\sou{mielo}. Substanco siropatra, sukroza, quan la abeli elaboras ek}
\l{\cel{3}la materii quin li rekoliis en la flori, e quin li igas ek-}
\l{\cel{3}fluar aden la alveoli di lia abeluyo por la nutro di su e di}
\l{\cel{3}la larvi. - FIS.}
\l{}
\l{mielato. Exuduro viskoza e sukroza, quan sekrecas, dum la perio-}
\l{\cel{3}di di aero-sikeso e di tero-sikeso, la folii di ula arbori.}
\l{}
\l{mieno. Aspekto di la vizajo. - DEFR.}
\l{}
\l{migrar. (\sou{netrans., ad, de)}. Diplasar su, livante ula lando por}
\l{\cel{3}instalar su en altra. - Su-diplasar di ula sorti de animali}
\l{\cel{3}qui iras ad altra klimati, sive periodale, segun la sezoni,}
\l{\cel{3}sive pro cirkonstanci acidentala. - EFIS.}
\l{}
\l{migreno. Doloro qua afektas, ordinare, nur parto di la kapo,}
\l{\cel{3}partikulare la regiono temporala, ed akompanesas multa-kaze}
\l{\cel{3}da perturbesi stomakala. - DEFIR.}
\l{}
\l{mikao. (\sou{mineral.}) Silikato aluminoza, kun metal-brilo, quan}
\l{\cel{3}onu povas separar facile ad lami diafana, per klivo.- DEFIS.}
\l{}
\l{mikologo.\cel{2}Autoro di traktato pri la fungi (champinioni).}
\l{}
\l{mikologio. Ta parto di botaniko qua studias la fungi (champi-}
\l{\cel{3}nioni).}
\l{}
\l{mikorizo. Fungo quan onu renkontras generale asociita kun radi-}
\l{\cel{3}ki di vejetanti (precipue la radiki di la vejetanti montala:}
\l{\cel{3}kastanieri, avelanieri, fagi, morusieri, monto-pini, e c.) -}
\l{\cel{3}DEFIS.}
}
